\documentclass[10pt,a4paper]{amsart}
\usepackage[margin=19mm]{geometry}
\usepackage{amsmath,amssymb,mathtools}
\usepackage[scr=rsfs,cal=euler]{mathalfa}
\usepackage{xfrac}
\usepackage[unicode,hidelinks,bookmarks=false]{hyperref}
\newcommand{\ie}{i.e.\ }
\newcommand{\cf}{cf.\ }
\newcommand{\resp}{respectively\ }
\newcommand{\Subsection}{\subsection}
\providecommand{\nobreakdash}{\nobreak\hbox{-}}
\setlength{\emergencystretch}{2em}
\multlinegap0pt
\usepackage{amssymb}
\usepackage{mathtools}
\usepackage{xcolor}
\newtheorem{definition}{Definition}
\newtheorem{lemma}{Lemma}
\newtheorem{proposition}{Proposition}
\newcommand{\E}{\mathbb{E}}
\newcommand{\KL}{\mathrm{KL}}
\newcommand{\KLinf}{\mathrm{KL}_{\mathrm{inf}}}
\newcommand{\Pcal}{\mathcal{P}}
\newcommand{\Tdet}{\mathsf T}
\newcommand{\X}{\mathsf{X}}
\newcommand{\bbF}{\mathbb F}
\newcommand{\bt}{\mathbf t}
\title{When the optimal betting wealth growth rate equals $\KLinf$}
\author{Ashwin Ram, Aaditya Ramdas}
\date{Source: arXiv:2604.25280v1 (CC BY 4.0)}
\begin{document}
\maketitle

\noindent\emph{Source and licence.} When the optimal betting wealth growth rate equals $\KLinf$. Version arXiv:2604.25280v1 (CC BY 4.0), distributed by arXiv under the Creative Commons Attribution 4.0 International licence, http://creativecommons.org/licenses/by/4.0/, as recorded in the arXiv licence metadata for this exact version and shown on the abstract page https://arxiv.org/abs/2604.25280v1. Full licence notice: \texttt{licenses/arxiv-2604.25280-CC-BY-4.0.txt}.\par\smallskip

\noindent\textit{Editorial scope.} Counted unit: the article's proposition prop:eprocess-rate (source0.tex lines 1318--1328). The article's complete original proof of it is delivered with it (source0.tex lines 1347--1417, one \texttt{proof} environment of 71 lines, which the article places away from the statement). No mathematics is written, repaired or re-derived: the statement and the proof are the article's own lines, in the article's own notation, and no retained line is substituted or re-flowed.  Retained uncounted as the source-local context the proof needs: lines 255--260; lines 653--658; lines 1332--1334.  Nothing was omitted from any retained span, so the package carries no omission record.  No retained line contains a citation, so the wrapper carries no bibliography and no bibliography entry.  No retained line contains a figure environment, an image-inclusion command or drawing code, so no image is delivered and the package declares no asset dependency.  Editorial formatting only, in addition: an AMS \texttt{amsart} preamble that restates the article's own declarations and macros used by the retained lines, namely the environments \texttt{definition}, \texttt{lemma}, \texttt{proposition} on separate counters, together with the standard packages those lines need.  The retained environments print in this document as Definition 1, Lemma 1, Proposition 1 in order of appearance, whereas the article prints the same items with the numbering of its own full document.  The complete native source of the article as distributed in the arXiv e-print for this exact version is bundled unchanged as \texttt{sources/199293/source0.tex}.
\begin{definition}[Blockwise test supermartingale]\label{def:test-supermartingale}
Let $\bt=(t_k)_{k\ge 0}\in \Tdet$. A blockwise test supermartingale for $\Pcal$ on $\bbF_{\bt}$ is a nonnegative process $(S_{t_k})_{k\ge 0}$ such that $S_{t_0}=1$, each $S_{t_k}$ is $\mathcal F_{t_k}$ measurable and for every $P\in\Pcal$ we have that for all $k\ge 1$ almost surely under $P^{\infty}$,
\[
\E_{P^{\infty}}[S_{t_k}\mid \mathcal F_{t_{k-1}}]\le S_{t_{k-1}}.
\]
\end{definition}

\begin{lemma}\label{lem:convex-empirical}
Let $P\in\mathcal M_1(\X)$ and let $C\subseteq\mathcal M_1(\X)$ be convex and weakly closed. Define $A_n:=\{\widehat Q_n\in C\}\subseteq \X^n$. Then it follows that for every $n\ge 1$,
\[
P^n(A_n)\ \le\ \exp\Bigl(-n\inf_{R\in C}\KL(R\|P)\Bigr).
\]
\end{lemma}

\begin{lemma}\label{lem:entropy-ineq}
Let $(\Omega, \mathcal F)$ be measurable and $\mu, \nu\in\mathcal M_1(\Omega)$. Now, then it follows that for any measurable $f$ with $\int e^{f}d\nu<\infty$ we have that $\int fd\mu \le \KL(\mu\|\nu)+\log\int e^{f}d\nu$.    
\end{lemma}

\begin{proposition}\label{prop:eprocess-rate}
Let $Q\in\mathcal M_1(\X)$ and let $\Pcal\subseteq\mathcal M_1(\X)$ be nonempty. Assume that $\KLinf(Q,\Pcal)<\infty$ and that $\Pcal$ is $\KLinf$ lower semicontinuous at $Q$. Let $d:=\KLinf(Q,\Pcal)=\inf_{P\in\Pcal}\KL(Q\|P)\in[0,\infty)$. Then, it follows that there exists a test supermartingale $(E_{t_k})_{k\ge 0}$ for $\Pcal$ and an associated nondecreasing sequences of sample sizes $(t_k)_{k\ge 0}$ with $t_0=0$ and $t_k\to\infty$ such that under $Q^\infty$,
\[
\lim_{k\to\infty}\frac{1}{t_k}\log E_{t_k} = d\quad{\text{almost surely.}}
\]
And also,
\[
\lim_{k\to\infty}\frac{1}{t_k}\E_{Q^\infty}[\log E_{t_k}]=d.
\]
If $d=0$ then we can take $E_{t_k}\equiv1$ for all $k$.
\end{proposition}

\begin{proof}[Proof of Proposition~\ref{prop:eprocess-rate}]
If $d=0$, then it follows that we can take $E_{t_k}\equiv 1$ and any $t_k\to\infty$, hence the conclusion of our proposition is immediate. Hence going forward, wlog we can assume that $d>0$. To this end, let us first begin by letting $(\varepsilon_k)_{k\ge 1}$ be a positive sequence such that $\varepsilon_k\downarrow 0$ and
\[
\frac{\sum_{j=1}^k\varepsilon_j m_j}{\sum_{j=1}^k m_j}\longrightarrow 0,
\]
for any sequence $(m_j)$ with $m_{j+1}\ge 2m_j$. To be concrete, one can take $\varepsilon_k:=\min\{2^{-k}, d/4\}$ so that $\varepsilon_k\downarrow 0$ and $d-2\varepsilon_k\ge d/2>0$ for all $k$. Then, let us define $\Phi(R):=\inf_{P\in\Pcal}\KL(R\|P)$. By $\KLinf$ lower semicontinuity of $\Phi$ at $Q$ that for each $k\ge 1$, the set $O_k:=\{R\in\mathcal M_1(\X):\Phi(R)>d-\varepsilon_k\}$ is a weakly open neighborhood of $Q$. Let us choose $\delta_k>0$ such that the closed and bounded Lipschitz ball
\[
C_k:=\overline B^{\mathrm{BL}}_{\delta_k}(Q)=\{R\in\mathcal M_1(\X):d_{\mathrm{BL}}(R,Q)\le\delta_k\}
\]
satisfies $C_k\subseteq O_k$. It follows then that $C_k$ is weakly closed and convex, and we have 
\[
\inf_{R\in C_k}\Phi(R)\ge d-\varepsilon_k.
\]
Since, $d_{\mathrm{BL}}(\widehat Q_m, Q)\to 0$ almost surely under $Q^{\infty}$, as $m\to\infty$, we have that $Q^m(\widehat Q_m\in C_k)\to 1$ for each $k$. Thanks to Lemma~\ref{lem:convex-empirical}, then
\[
\sup_{P\in\Pcal} P^m(\widehat Q_m\in C_k) \le \exp\Bigl(-m\inf_{R\in C_k}\Phi(R)\Bigr) \le \exp\bigl(-m(d-\varepsilon_k)\bigr)
\]
Let us now choose the block lengths $m_k$ recursively so that $m_k\ge 2m_{k-1}$ for $k\ge 2$ and
\begin{equation}\label{eq:eproc-Qgood}
Q^{m_k}\bigl(\widehat Q_{m_k}\in C_k\bigr)\ge 1-2^{-k}    
\end{equation}
Now define $A_k:=\{\widehat Q_{m_k}\in C_k\}\subseteq X^{m_k}$. Then
\begin{equation}\label{eq:eproc-betabound}
\sup_{P\in\Pcal}P^{m_k}(A_k) \le \exp\bigl(-m_k(d-\varepsilon_k)\bigr) \le \exp\bigl(-m_k(d-2\varepsilon_k)\bigr)   
\end{equation}
Now, let us define the cumulative sample sizes $t_0:=0$ and $t_k:=\sum_{j=1}^k m_j$. Now let $(\gamma_k)_{k\ge 1}\subset(0,1)$ be summable. Define the $k$th block factor by
\begin{equation}\label{eq:eproc-blockE}
E^{(k)}(x^{m_k}):=(1-\gamma_k)\exp\bigl(m_k(d-2\varepsilon_k)\bigr)\mathbf 1\!\{\widehat Q_{m_k}(x^{m_k})\in C_k\} +\gamma_k    
\end{equation}

Clearly, this is strictly positive and depends only on the data inside the $k$th block. With this now, let us place the observations into successive disjoint blocks. Meaning, block $k$ consists of coordinates $X_{t_{k-1}+1},\dots,X_{t_k}$. Define $E_{t_0}:=1$ and for $k\ge 1$,
\[
E_{t_k}:=\prod_{j=1}^k E^{(j)}\bigl(X_{t_{j-1}+1},\dots,X_{t_j}\bigr).
\]
Then $(E_{t_k})_{k\ge 0}$ is adapted to the $(\mathcal F_{t_k})_{k\ge 0}$. Let us verify that the $e$-process property holds for every $P\in\Pcal$. Meaning, take some $P\in\Pcal$. Because $P^\infty$ is iid, the block $(X_{t_{k-1}+1},\dots, X_{t_k})$ is independent of $\mathcal F_{t_{k-1}}$ and has law $P^{m_k}$. Hence it follows that $\E_{P^\infty}[E_{t_k}\mid \mathcal F_{t_{k-1}}]= E_{t_{k-1}}\E_{P^{m_k}}\left[E^{(k)}\right]$. Therefore we can use the definition \eqref{eq:eproc-blockE} and the bound \eqref{eq:eproc-betabound} to give us that
\begin{align*}
\E_{P^{m_k}}\left[E^{(k)}\right]&=(1-\gamma_k)\exp\bigl(m_k(d-2\varepsilon_k)\bigr)P^{m_k}(A_k)+\gamma_k\\
&\le(1-\gamma_k)\exp\bigl(m_k(d-2\varepsilon_k)\bigr)\sup_{P'\in\Pcal}P'^{m_k}(A_k)+\gamma_k\\
&\le(1-\gamma_k)\cdot 1+\gamma_k=1.
\end{align*}
Then, $\E_{P^\infty}[E_{t_k}\mid\mathcal F_{t_{k-1}}]\le E_{t_{k-1}}$ almost surely. Hence $(E_{t_k})_{k\ge 0}$ is a blockwise test supermartingale as in Definition~\ref{def:test-supermartingale}. We now need to show the almost-sure growth rate under $Q^\infty$. Under $Q^\infty$, the blocks are i.i.d.\ with law $Q^{m_k}$. They are independent across $k$ but not identical since $k$ can absolutely vary. Now let $B_k$ be the event that the $k$th block belongs to $A_k$. That is,
\[
B_k:=\Bigl\{\widehat Q_{m_k}\bigl(X_{t_{k-1}+1},\dots,X_{t_k}\bigr)\in U_k\Bigr\}.
\]
By \eqref{eq:eproc-Qgood} it follows that $Q^\infty(B_k^c)\le 2^{-k}$. Henceby Borel Cantelli Lemma 1, $Q^\infty(B_k\text{ eventually})=1$. On the event that $B_k$ holds for all $k\ge K(\omega)$, for such $k$ we have that $E^{(k)}=(1-\gamma_k)\exp\bigl(m_k(d-2\varepsilon_k)\bigr)+\gamma_k$. Therefore,
\[
\log E^{(k)} \in \Bigl[m_k(d-2\varepsilon_k)+\log(1-\gamma_k),\;m_k(d-2\varepsilon_k)\Bigr],
\]
since $\log\bigl((1-\gamma_k)e^{a}+\gamma_k\bigr)\le \log(e^a)=a$, and $\log\bigl((1-\gamma_k)e^{a}+\gamma_k\bigr)\ge \log\bigl((1-\gamma_k)e^{a}\bigr)=a+\log(1-\gamma_k)$. Let us now sum over all $k\le n$ and then divide by $t_n=\sum_{j=1}^n m_j$. If we do this we get that for all large $n$ on this event that
\[
d-2\frac{\sum_{k=1}^n \varepsilon_k m_k}{t_n} +\frac{\sum_{k=1}^n \log(1-\gamma_k)}{t_n} \le \frac{1}{t_n}\log E_{t_n} \le d-2\frac{\sum_{k=1}^n \varepsilon_k m_k}{t_n} +\frac{C(\omega)}{t_n}.
\]
Note that here $C(\omega)$ is simply some finite constant which accounts for the fact that there are finitely many exceptional blocks $k<K(\omega)$. Since $t_n\to\infty$, we know that $C(\omega)/t_n\to 0$. Further, $\sum_k |\log(1-\gamma_k)|<\infty$ because $\sum_k \gamma_k<\infty$. Hence, $\sum_{k=1}^n \log(1-\gamma_k)=O(1)$ and thus $\frac{1}{t_n}\sum_{k=1}^n \log(1-\gamma_k)\to 0$. Lastly, by our choice of $(\varepsilon_k)$ in conjunction with the geometric growth of $(m_k)$ we have that
\[
\frac{\sum_{k=1}^n \varepsilon_k m_k}{t_n}\longrightarrow 0.
\]
If we take all that we have shown now together we have $\frac{1}{t_n}\log E_{t_n}\to d$ almost surely under $Q^\infty$. It remains to now show the expected log-growth rate under $Q^\infty$. Because we always have $\log E_{t_n}=\sum_{k=1}^n \log E^{(k)}$, it follows that by linearity, $\E_{Q^\infty}[\log E_{t_n}]=\sum_{k=1}^n \E_{Q^{m_k}}[\log E^{(k)}]$. On $B_k$ we have that $\log E^{(k)}\ge m_k(d-2\varepsilon_k)+\log(1-\gamma_k)$. And on $B_k^c$, $\log E^{(k)}=\log\gamma_k$. Therefore we get that
\[
\E_{Q^{m_k}}[\log E^{(k)}] \ge Q^{m_k}(A_k)\Bigl(m_k(d-2\varepsilon_k)+\log(1-\gamma_k)\Bigr) +Q^{m_k}(A_k^c)\log\gamma_k.
\]
Recall that from \eqref{eq:eproc-Qgood}, $Q^{m_k}(A_k)\ge 1-2^{-k}$. Let us also use the two facts that $\log(1-\gamma_k)=O(\gamma_k)$ with $\sum_k\gamma_k<\infty$, and $\sum_k 2^{-k}|\log\gamma_k|<\infty$. Hence,
\[
\frac{1}{t_n}\E_{Q^\infty}[\log E_{t_n}] \ge d-2\frac{\sum_{k=1}^n \varepsilon_k m_k}{t_n} - o(1).
\]
Separately, $(E_{t_k})$ is a blockwise $e$-process for $\Pcal$. Hence it follows that each $E_{t_k}$ is a block $e$-variable for the null $\Pcal$ at block time $k$. Thus for every $P\in\Pcal$ we have that $\E_{P^\infty}[E_{t_k}]\le 1$. Hence by Lemma~\ref{lem:entropy-ineq} applying $\mu=Q^{t_k}$, $\nu=P^{t_k}$ and $f=\log E_{t_k}$ (so $e^f=E_{t_k}$) we get that
\[
\E_{Q^\infty}[\log E_{t_k}]\le \KL(Q^{t_k}\|P^{t_k})+\log \E_{P^\infty}[E_{t_k}].
\]
Necessarily $(E_{t_k})$ is an $e$-process for $\Pcal$, which means that $\E_{P^\infty}[E_{t_k}]\le 1$. Therefore, $\E_{Q^\infty}[\log E_{t_k}]\le \KL(Q^{t_k}\|P^{t_k})$. Now if we take the $\inf_{P\in\Pcal}$ and use tensorization, $\E_{Q^\infty}[\log E_{t_k}]\le t_k d$.
Dividing by $t_k$ and combining with the lower bound gives us equality; that is, $\lim_{k\to\infty}\frac{1}{t_k}\E_{Q^\infty}[\log E_{t_k}]=d$. This completes the proof.
\end{proof}

\end{document}
